\documentclass[10pt]{article}
\usepackage[utf8]{inputenc}
\usepackage[T1]{fontenc}
\usepackage{graphicx}
\usepackage[export]{adjustbox}
\graphicspath{ {./images/} }
\usepackage{caption}
\usepackage{amsmath}
\usepackage{amsfonts}
\usepackage{amssymb}
\usepackage[version=4]{mhchem}
\usepackage{stmaryrd}
\usepackage{hyperref}
\hypersetup{colorlinks=true, linkcolor=blue, filecolor=magenta, urlcolor=cyan,}
\urlstyle{same}

\begin{document}
\captionsetup{singlelinecheck=false}
Introduction

\section*{Forward Induction}
\begin{itemize}
  \item Subgame Perfection and Sequential Equilibrium are ˋˋbackward induction" concepts, i.e., players behavior towards the end of the game determines players behavior earlier in the game.
  \item Backward induction ideas have been formalized in various notions including dynamic programming and sequential rationality.
  \item Kohlberg and Mertens (1986) proposed a notion of "forward induction:" a player's options early in a game may determine (and help predict) the player's behavior later in the game. Their refinement is termed Stability.
\end{itemize}

\section*{Various expressions of forward induction}
\begin{itemize}
  \item Iterative Elimination of Weakly Dominated Strategies (IEWDS).
  \item Equilibrium Dominance.
  \item The Intuitive Criterion, Cho and Kreps (1987).
  \item Divinity, Banks and Sobel (1987).
  \item Never a Weak Best Response Criterion (NWBR).
  \item Various formulations of stability, Kohlberg and Mertens (1986)\}, \textbackslash textcite\{govindan-wilson2009.
\end{itemize}

\section*{Example}
\begin{figure}[h]
\begin{center}
  \includegraphics[alt={},max width=\textwidth]{a88a03fe-5705-4253-8e40-af3ffb35fb27-05_534_1431_343_270}
\captionsetup{labelformat=empty}
\caption{Figure 1}
\end{center}
\end{figure}

\begin{itemize}
  \item Three strategy profiles that are part of sequential equilibria. To prove these are SE, beliefs must be provided to construct consistent assessments.
\end{itemize}

\begin{enumerate}
  \item $[(A, Y), w]$; 2. $[(A,(.8, .2)),(3 / 8,5 / 8)]$; and 3. $[(B, X), z]$.
\end{enumerate}

\section*{Example Figure 1}
\begin{enumerate}
  \item $[(A, Y), w]$
\end{enumerate}

$$
\begin{aligned}
\sigma_{1}\left(\cdot \mid h_{1}\right) & =\left(\sigma_{1}\left(A \mid h_{1}\right), \sigma_{1}\left(B \mid h_{1}\right)\right)=(1,0) \\
\sigma_{1}\left(\cdot \mid h_{1}^{\prime}\right) & =\left(\sigma_{1}\left(X \mid h_{1}\right), \sigma_{1}\left(Y \mid h_{1}\right)\right)=(0,1) \\
\sigma_{2}\left(\cdot \mid h_{2}\right) & =\left(\sigma_{2}\left(Z \mid h_{2}\right), \sigma_{2}\left(W \mid h_{2}\right)\right)=(0,1) \\
\mu_{1}\left(h_{1} \mid h_{1}\right) & =1 \text { and } \mu_{1}\left(h_{1}^{\prime} \mid h_{1}^{\prime}\right)=1 \\
\sigma_{1}\left(\cdot \mid h_{1}\right) & =\left(\sigma_{1}\left(A \mid h_{1}\right), \sigma_{1}\left(B \mid h_{1}\right)\right)=(1,0) \\
\mu_{2}\left((B, X) \mid h_{2}\right) & =0 \text { and } \mu_{2}\left((B, Y) \mid h_{2}\right)=1 \text { and }
\end{aligned}
$$

$$
\begin{aligned}
\sigma_{1}^{n}\left(\cdot \mid h_{1}\right) & =\left(1-\frac{1}{n}, \frac{1}{n}\right) \\
\sigma_{1}^{n}\left(\cdot \mid h_{1}^{\prime}\right) & =\left(\frac{1}{n}, 1-\frac{1}{n}\right) \\
\sigma_{2}^{n}\left(\cdot \mid h_{2}\right) & =\left(\frac{1}{n}, 1-\frac{1}{n}\right) \\
\mu_{2}^{n}\left(\cdot \mid h_{2}\right) & =\left(\frac{1}{n}, 1-\frac{1}{n}\right)
\end{aligned}
$$

because $\mu_{2}^{n}\left((B, X) \mid h_{2}\right)=\frac{1 / n^{2}}{1 / n}=1 / n$.\\
2. $[(A,(.8, .2)),(3 / 8,5 / 8)]$.

$$
\begin{array}{lr}
\sigma_{1} \longrightarrow & (1,0),(0.8,0.2) \\
\sigma_{2} \longrightarrow & (3 / 8,5 / 8)
\end{array}
$$

\begin{enumerate}
  \setcounter{enumi}{2}
  \item $[(B, X), z]$.
\end{enumerate}

$$
\begin{array}{lr}
\sigma_{1} \longrightarrow & (0,1),(1,0) \\
\sigma_{2} \longrightarrow & (1,0)
\end{array}
$$

\section*{Forward induction. Intuition}
\begin{itemize}
  \item Forward induction arguments eliminate the first two equilibria.
  \item Deviations from proposed equilibrium reveal the intentions of the deviating player in the continuation game.
  \item In the first equilibrium, player 1's equilibrium payoff is 4 . Suppose player 1 deviates from this equilibrium and chooses $B$. What can player 1 expect?
\end{itemize}

\section*{Example}
\begin{itemize}
  \item Iterative elimination of weakly dominated strategies (IEWDS)
\end{itemize}

\begin{figure}[h]
\begin{center}
  \includegraphics[alt={},max width=\textwidth]{a88a03fe-5705-4253-8e40-af3ffb35fb27-10_526_1419_565_278}
\captionsetup{labelformat=empty}
\caption{Figure 2}
\end{center}
\end{figure}

\begin{center}
\begin{tabular}{lll}
 & $z$ & $w$ \\
\hline
$A, X$ & $(4,0)$ & $(4,0)$ \\
\hline
$A, Y$ & $(4,0)$ & $(4,0)$ \\
\hline
$B, X$ & $(5,1)$ & $(0,0)$ \\
\hline
$B, Y$ & $(0,0)$ & $(3,4)$ \\
\hline
\end{tabular}
\end{center}

\begin{itemize}
  \item 1's strategy ( $B, Y$ ) is dominated by ( $A, Y$ ).
  \item Then 2's strategy $W$ becomes weakly dominated by $Z$.
  \item The surviving equilibrium is $[(B, X), Z]$.
\end{itemize}

\section*{Example: Equilibrium dominance}
\begin{itemize}
  \item Is $[(A Y), w]$ intuitive? Its payoffs are $(4,0)$.
  \item Subgame starting at $h_{1}^{\prime}$ has 3 NE:
\end{itemize}

\begin{center}
\begin{tabular}{llll}
 & $(X, Y)$ & $(z, w)$ & Player 1's Payoff \\
\hline
a. & $(0,1)$ & $(0,1)$ & 3 \\
\hline
b. & $(0.8,0.2)$ & $(3 / 8,5 / 8)$ & $15 / 8$ \\
\hline
c. & $(1,0)$ & $(1,0)$ & 5 \\
\hline
\end{tabular}
\end{center}

\begin{itemize}
  \item Only option c improves player 1's payoff since $5>4$.
  \item Player 1 would only deviate to $B$ (from $A$ ) only if she intends to benefit. That suggests that 1 would continue with $X$.
  \item At $h_{2}$, player 2 ought to believe that player 1 chose $(B, X)$.\\
-If 2 reasons as described, 1 should choose $(B, X)$ (i.e., $(0,1)(1,0)$ ). Player 2 should respond $(1,0)$.
  \item Player 2 infers 1 's behavior from the fact that the player deciding at $h_{1}^{\prime}$ is the same as the player deciding at $h_{1}$.
  \item This suggests that an analytic formulation must depend on the Normal form and not on the Agent Normal Form.
  \item In the example, the intuition behind "forward induction" may be formalized by the iterative elimination of weakly dominated strategies.
  \item In other games, IEWDS yields a backward induction equilibrium.
  \item Still in other games, IEWDS yields either a backward induction equilibrium or a forward induction equilibrium depending on the order of elimination of strategies.
  \item Forward induction is based on the principle that deviations should be explained rationally; backward induction on the principle that deviations are mistakes.
\end{itemize}

\section*{Example: forward-backward conflict}
\begin{itemize}
  \item Conflict between forward and backward induction IEWDS.
\end{itemize}

Myerson (1991), pp. 191-196.

\begin{figure}[h]
\begin{center}
  \includegraphics[alt={},max width=\textwidth]{a88a03fe-5705-4253-8e40-af3ffb35fb27-15_522_1596_694_270}
\captionsetup{labelformat=empty}
\caption{Figure 3}
\end{center}
\end{figure}

\begin{enumerate}
  \item Process starts with 2 , strategy $(E, z)$. Only one strategy is eliminated in each iteration:
\end{enumerate}

\begin{itemize}
  \item 2's strategy $(E, z)$ weakly dominated by all others,
  \item 1's strategy ( $B, X$ ) weakly dominated by all others, $(B, Y)$ also dominated but it is not eliminated at this stage.
  \item 2's strategy ( $D, \cdot$ ) is weakly dominated by ( $E, w$ ),
  \item OUTCOME: $[(A, \cdot),(E, w)]$. This is the backward induction equilibrium.
\end{itemize}

\begin{enumerate}
  \setcounter{enumi}{1}
  \item The process starts with player 1 , strategy $(B, Y)$.
\end{enumerate}

\begin{itemize}
  \item 1's strategy $(B, Y)$ weakly dominated by $(A, \cdot)$,
  \item 2's strategy $(E, z)$ weakly dominated by all others,
  \item 1's strategy ( $B, X$ ) weakly dominated by ( $A$, ),
  \item OUTCOMES: $[(A, \cdot)(D, \cdot) ;(A, \cdot)(E, w)]$.
\end{itemize}

The forward induction equilibrium is contained in the set of outcomes obtained by IEWDS.

\section*{References}
Banks, Jeffrey S., and Joel Sobel. 1987. "Equilibrium Selection in Signaling Games." Econometrica 55 (3): 647-61. \href{https://doi.org/10.2307/1913604}{https://doi.org/10.2307/1913604}.

Cho, In-Koo, and David Kreps. 1987. "Signaling Games and Stable Equilibria." Quarterly Journal of Economics 102: 179-222.\\
Kohlberg, Elon, and Jean-François Mertens. 1986. "On the Strategic Stability of Equilibria." Econometrica 54 (5): 1003-37.


\end{document}